\section*{第 \ref{ch:b3:groups} 章——群论}

\begin{solution}{exo:b3:groups:1}
(a) 设 $[G:H] = 2$。对 $g \in H$，$gH = H = Hg$。对 $g \notin
H$：两个左陪集为 $H$ 与 $gH$，于是 $gH = G \setminus H$；
类似地，$Hg = G \setminus H$。因此 $gH = Hg$ 对所有 $g$ 成立，即 $H
\trianglelefteq G$。

(b) 由 (a)，$G/H$ 是 $2$ 阶群；元素 $\bar x$
满足 $\bar x^2 = \bar e$，即 $x^2 \in H$。

(c) 设 $H \leq A_4$ 且 $\abs H = 6$，从而指数为 $2$。由
(b)，$\sigma^2 \in H$ 对每个 $\sigma \in A_4$ 成立。每个 $3$-轮换
都是这样的平方：若 $\sigma^3 = e$，则 $\sigma =
\sigma^4 = (\sigma^2)^2$。于是 $H$ 包含全部八个 $3$-轮换，
它们都在 $A_4$ 中：$\abs H \geq 8 > 6$，矛盾。（拉格朗日定理的逆
命题在首个可能之处即告失效：$6 \mid 12$。）
\end{solution}

\begin{solution}{exo:b3:groups:2}
设 $G/Z(G) = \langle gZ(G) \rangle$。于是每个 $x \in G$ 可写为
$x = g^k z$，其中 $k \in \Z$，$z \in Z(G)$。对 $x = g^kz$，$y =
g^l z'$：
\[
xy = g^k z\, g^l z' = g^{k+l} z z' = g^l z'\, g^k z = yx,
\]
中心元与一切元素交换：$G$ 是阿贝尔群。

$p^2$ 阶群：$Z(G) \neq \{e\}$（由\cref{thm:b3:groups:pfixed}），故
$\abs{Z(G)} \in \{p, p^2\}$。若其阶为 $p$，则 $G/Z(G)$
为 $p$ 阶群，从而为循环群，迫使 $G$ 为阿贝尔群且 $Z(G) =
G$，阶为 $p^2$——矛盾。故 $Z(G) = G$。

非阿贝尔的 $p^3$ 阶群：\emph{海森堡群}
\[
H_p = \left\{ \begin{pmatrix} 1 & a & c\\ 0 & 1 & b\\ 0 & 0 & 1
\end{pmatrix} : a, b, c \in \mathbb F_p \right\}
\leq GL_3(\mathbb F_p),
\]
其阶为 $p^3$（$a, b, c$ 可自由选取；封闭性与逆元
由直接计算得出）。它是非阿贝尔的：两个初等矩阵
$I + E_{12}$ 与 $I + E_{23}$ 的\omterm{def:b3:groups:derived}{换位子}为 $I + E_{13} \neq I$。
\end{solution}

\begin{solution}{exo:b3:groups:3}
(a) 同态 $f \colon \Z/n\Z \to \Z/n\Z$ 由 $k =
f(\bar 1)$（此时 $f(\bar m) = m\bar k$），且每个 $\bar k$ 都确定
一个同态。它是双射当且仅当 $\bar k$ 生成 $\Z/n\Z$，当且仅当
$\gcd(k,n) = 1$，当且仅当 $\bar k \in (\Z/n\Z)^\times$。复合
对应于乘法：$f_k \circ f_l = f_{kl}$。因此
$\operatorname{Aut}(\Z/n\Z) \cong (\Z/n\Z)^\times$。

(b) 映射 $\iota\colon G \to \operatorname{Aut}(G)$，$g \mapsto
\iota_g$，是同态：$\iota_g \circ \iota_h = \iota_{gh}$。其
像为 $\operatorname{Inn}(G)$；其核为 $\{g : gxg^{-1} =
x\ \forall x\} = Z(G)$。由第一同构定理得
$\operatorname{Inn}(G) \cong G/Z(G)$。关于
$\operatorname{Aut}(G)$ 中的正规性：对 $\alpha \in \operatorname{Aut}(G)$，
\[
(\alpha \circ \iota_g \circ \alpha^{-1})(x)
= \alpha\bigl(g\,\alpha^{-1}(x)\,g^{-1}\bigr)
= \alpha(g)\, x\, \alpha(g)^{-1} = \iota_{\alpha(g)}(x).
\]
\end{solution}

\begin{solution}{exo:b3:groups:4}
(a) 考虑 $\mu \colon H \times K \to HK$，$(h,k) \mapsto hk$，
由定义，它是满射。固定 $h_0k_0 \in HK$，则 $hk = h_0k_0
\iff h_0^{-1}h = k_0 k^{-1} \in H \cap K$。记 $u =
h_0^{-1}h$，则 $h_0k_0$ 的纤维为 $\{(h_0u,\, u^{-1}k_0) : u
\in H \cap K\}$，其元素个数为 $\abs{H \cap K}$。因此 $\abs
H\,\abs K = \abs{H\times K} = \abs{HK}\,\abs{H \cap K}$。

(b) 若 $HK$ 为子群，则 $KH \subseteq HK$，因为 $kh =
\bigl(h^{-1}k^{-1}\bigr)^{-1} \in (HK)^{-1} = HK$；又 $HK
\subseteq KH$，取逆可得：$HK = (HK)^{-1} \subseteq
(KH)^{-1}\dots$ 更直接地，对 $hk \in HK$，$(hk)^{-1} =
k^{-1}h^{-1} \in KH$，故 $HK = (HK)^{-1} \subseteq KH$；两个
包含关系给出 $HK = KH$。反之，若 $HK = KH$，则封闭性：
$(hk)(h'k') = h(kh')k' \in h(HK)k' = (hH)(Kk') \subseteq HK$；
逆元：$(hk)^{-1} = k^{-1}h^{-1} \in KH = HK$；且 $e \in HK$，
故为子群。若 $K \trianglelefteq G$，则 $hK = Kh$ 对所有
$h$，有 $HK = KH$ 自动成立。

(c) 对 $h \in H$，$k \in K$，\omterm{def:b3:groups:derived}{换位子} $[h,k] = hkh^{-1}
k^{-1}$ 等于 $(hkh^{-1})k^{-1} \in K$（$K$ 正规）并且
$h(kh^{-1}k^{-1}) \in H$（$H$ 正规），从而位于 $H \cap K =
\{e\}$ 中：$hk = kh$。
\end{solution}

\begin{solution}{exo:b3:groups:5}
用两种方式计算 $E = \{(g,x) \in G \times X : g \cdot x = x\}$ 的元素个数：
\[
\abs E = \sum_{g \in G} \abs{\operatorname{Fix}(g)}
= \sum_{x \in X} \abs{G_x}
= \sum_{x \in X} \frac{\abs G}{\abs{\mathcal O_x}}
= \abs G \sum_{\mathcal O \text{ 轨道}} \sum_{x \in \mathcal O}
\frac{1}{\abs{\mathcal O}}
= \abs G \cdot \#\{\text{轨道}\},
\]
利用轨道--稳定子定理（$\abs{G_x} = \abs G/\abs{\mathcal O_x}$）
以及\omterm{def:b3:groups:action}{轨道}的划分。

项链计数：设 $X$ 为映射 $\Z/p\Z \to
\{1, \dots, a\}$（即 $p$ 个位置的染色），$\abs X = a^p$，
$\Z/p\Z$ 通过旋转作用。恒等元固定全部 $a^p$
种染色。非零旋转 $\bar k \neq \bar 0$ 生成 $\Z/p\Z$（$p$
为质数），故它固定的染色对\emph{全部}旋转都不变，
从而为常值染色，被固定的染色恰有 $a$ 种。由伯恩赛德引理：
\[
\#\{\text{轨道}\} = \frac{a^p + (p-1)a}{p} \in \N ,
\]
故 $p \mid a^p + (p-1)a$，即 $p \mid a^p - a$：费马
小定理，仅凭计数即得。
\end{solution}

\begin{solution}{exo:b3:groups:6}
群阶 $15 = 3 \cdot 5$：$n_3 \mid 5$ 且 $n_3 \equiv 1 \pmod 3$
推出 $n_3 = 1$；$n_5 \mid 3$ 且 $n_5 \equiv 1 \pmod 5$ 推出
$n_5 = 1$。\omterm{def:b3:groups:sylow}{西罗子群} $P_3, P_5$ 均正规，且两者交集平凡
（阶互素），且 $\abs{P_3P_5} = 15$
（见\cref{exo:b3:groups:4}(a)）。由\cref{prop:b3:groups:direct}，$G
\cong \Z/3\Z \times \Z/5\Z \cong \Z/15\Z$（中国剩余定理）。

群阶 $45 = 3^2 \cdot 5$：$n_3 \mid 5$，$n_3 \equiv 1 \pmod 3$
给出 $n_3 = 1$；$n_5 \mid 9$，$n_5 \equiv 1 \pmod 5$ 给出 $n_5 =
1$。故 $G \cong P_3 \times P_5$，其中 $\abs{P_3} = 9 = 3^2$ 且
$\abs{P_5} = 5$：两者均为阿贝尔群（由\cref{thm:b3:groups:pfixed} 给出
$p^2$ 阶情形，素数阶群为循环群），从而 $G$ 亦然。
\end{solution}

\begin{solution}{exo:b3:groups:7}
群阶 $30 = 2 \cdot 3 \cdot 5$。$n_5 \mid 6$，$n_5 \equiv 1 \pmod
5$：$n_5 \in \{1, 6\}$；$n_3 \mid 10$，$n_3 \equiv 1 \pmod 3$：
$n_3 \in \{1, 10\}$。设 $G$ 为\omterm{def:b3:groups:simple}{单群}，则 $n_5 = 6$ 且 $n_3 =
10$。素数阶 $p$ 的两个不同子群交集平凡
（交集是 $\Z/p\Z$ 的真子群），故 6 个
西罗 $5$-子群共含 $6 \times 4 = 24$ 个 $5$ 阶元，
而 10 个西罗 $3$-子群共含 $10 \times 2 = 20$ 个
$3$ 阶元，两类元素合计 $24 + 20 = 44 > 30 - 1$ 个非单位元——
矛盾。故 $n_5 = 1$ 或 $n_3 = 1$，即存在正规\omterm{def:b3:groups:sylow}{西罗子群}。

群阶 $56 = 2^3 \cdot 7$。$n_7 \mid 8$，$n_7 \equiv 1 \pmod 7$：
$n_7 \in \{1, 8\}$。若 $n_7 = 8$，则西罗 $7$-子群共含 $8
\times 6 = 48$ 个 $7$ 阶元，恰余 $56 - 48 =
8$ 个其他元素。西罗 $2$-子群阶为 $8$，其元素
皆属于此集合，故它\emph{恰为}该集合：$n_2 = 1$。
故 $n_7 = 1$ 或 $n_2 = 1$，此群决不为\omterm{def:b3:groups:simple}{单群}。
\end{solution}

\begin{solution}{exo:b3:groups:8}
(a) 作用 $g \cdot xH = gxH$ 给出同态 $\rho \colon G
\to \mathfrak S(G/H) \cong S_n$。其核为
\[
\ker\rho = \{g : \forall x \in G,\ gxH = xH\}
= \{g : \forall x,\ x^{-1}gx \in H\}
= \bigcap_{x \in G} xHx^{-1},
\]
它是\omterm{def:b3:groups:normal}{正规子群}（即核），且含于 $H$（取 $x = e$）。若
$N \trianglelefteq G$ 且 $N \subseteq H$，则对每个 $x$ 有 $N
= xNx^{-1} \subseteq xHx^{-1}$，故 $N \subseteq \ker\rho$：核
即是此类子群中最大者。

(b) 设 $[G:H] = p$ 为整除 $\abs G$ 的最小素数，$K =
\ker\rho \subseteq H$。则 $G/K$ 嵌入 $S_p$，故 $[G:K]$
整除 $p!$。又 $[G:K] = [G:H]\,[H:K] = p\,[H:K]$，故 $[H:K]$
整除 $(p-1)!$。但 $[H:K]$ 整除 $\abs G$，其素因数
均 $\geq p$，而 $(p-1)!$ 的素因数均 $< p$，
故 $[H:K] = 1$，即 $H = K = \ker \rho$ 为
\omterm{def:b3:groups:normal}{正规子群}。

(c) 设作用可迁，$x \in X$。映射 $\Phi
\colon G/G_x \to X$，$gG_x \mapsto g \cdot x$，良定义且
双射（轨道--稳定子定理；\omterm{def:b3:groups:action}{轨道}即为 $X$），并且它
与作用相容：$\Phi(h \cdot gG_x) = \Phi(hgG_x) = (hg)
\cdot x = h \cdot \Phi(gG_x)$。
\end{solution}

\begin{solution}{exo:b3:groups:9}
(a) $D(G)$ 正规且\omterm{thm:b3:groups:quotient}{商群}为阿贝尔群
（见\cref{prop:b3:groups:derived}）；若 $N \trianglelefteq G$ 使得
$G/N$ 为阿贝尔群，则由同一命题得 $D(G) \subseteq N$：
$D(G)$ 即为最小者。泛性质：设 $f \colon G \to A$
其中 $A$ 为阿贝尔群。则 $f([x,y]) = [f(x), f(y)] = e$，故 $D(G)
\subseteq \ker f$，由\cref{thm:b3:groups:firstiso} 知 $f =
\bar f \circ \pi$ 经 $G^{\mathrm{ab}}$ 分解，且因 $\pi$
为满射，分解唯一。

(b) \omterm{def:b3:groups:derived}{换位子}为偶置换，故 $D(S_n) \subseteq A_n$。
反之，每个 $3$-轮换都是\omterm{def:b3:groups:derived}{换位子}：
\[
\bigl[(a\,b),\,(a\,c)\bigr] = (a\,b)(a\,c)(a\,b)(a\,c) =
(a\,b\,c),
\]
（对 $a, b, c$ 直接验证），且 $3$-轮换生成 $A_n$
（见\cref{lem:b3:groups:threecycles}），故 $D(S_n) = A_n$（$n \geq 3$），
且 $S_n^{\mathrm{ab}} \cong S_n/A_n \cong \Z/2\Z$。（$n = 2$ 时：
$S_2$ 为阿贝尔群，$D(S_2) = \{e\}$，$S_2^{\mathrm{ab}} = S_2
\cong \Z/2\Z$——公式 $S_n^{\mathrm{ab}} \cong \Z/2\Z$
对所有 $n \geq 2$ 成立。）

$Q_8$：\omterm{thm:b3:groups:quotient}{商群} $Q_8/\{\pm 1\}$ 的阶为 $4$，故为
阿贝尔群，从而 $D(Q_8) \subseteq \{\pm 1\}$；又 $[\mathrm i, \mathrm
j] = \mathrm i \mathrm j \mathrm i^{-1}\mathrm j^{-1} = \mathrm
i\mathrm j(-\mathrm i)(-\mathrm j) = (\mathrm i \mathrm j)^2 =
\mathrm k^2 = -1$，故 $D(Q_8) = \{\pm 1\}$，且
$Q_8^{\mathrm{ab}} \cong (\Z/2\Z)^2$（阶为 $4$，方次数为 $2$：
$\mathrm i, \mathrm j$ 所在的陪集平方均为 $\bar 1$）。
\end{solution}

\begin{solution}{exo:b3:groups:10}
对 $\abs G$ 作归纳；当 $\abs G = p$ 时，唯一的真子群为
$H = \{e\}$，且 $N_G(\{e\}) = G \supsetneq
\{e\}$。设 $Z = Z(G) \neq \{e\}$（见\cref{thm:b3:groups:pfixed}）。

\emph{若 $Z \not\subseteq H$}：取 $z \in Z \setminus H$；$z$
与 $H$ 可交换，故 $zHz^{-1} = H$，从而 $z \in N_G(H) \setminus
H$。

\emph{若 $Z \subseteq H$}：转到 $\bar G = G/Z$，它是阶更小的 $p$-群，
且 $\bar H = H/Z \subsetneq \bar G$（由对应定理
），由归纳假设 $N_{\bar G}(\bar H) \supsetneq \bar H$；
取 $\bar g \in N_{\bar G}(\bar H) \setminus \bar H$ 及其一个提升
$g$。则 $g \notin H$，且 $gHg^{-1} \subseteq HZ = H$：事实上
由 $\overline{ghg^{-1}} = \bar g \bar h \bar g^{-1} \in \bar H$ 知
$ghg^{-1} \in HZ = H$（因 $Z \subseteq H$）。故 $g \in
N_G(H)\setminus H$。

极大子群：若 $M$ 为极大子群，则 $N_G(M) \supsetneq M$
必有 $N_G(M) = G$，即 $M \trianglelefteq G$。于是 $G/M$ 为
一个 $p$-群，除平凡子群与自身外不含其他子群（由对应定理
与极大性）。取 $\bar x \neq \bar e$（属于 $G/M$），其阶为 $p^k$；
则 $\bar x^{p^{k-1}}$ 生成一个 $p$ 阶子群，
这必为全体，故 $\abs{G/M} = p$。
\end{solution}

\begin{solution}{exo:b3:groups:11}
(a) $\abs{A_5} = 60$。$A_5$ 中的轮换型有：$e$；双
对换，共 $\frac{1}{2}\binom{5}{1}\binom{4}{2}\cdot 1 = 15$
个（$5 \cdot 3$ 种方式：先选不动点，再两两配对）；
$3$-轮换，$\frac{5 \cdot 4 \cdot 3}{3} = 20$；$5$-轮换，$4! =
24$。

$S_5$ 的一个\omterm{ex:b3:groups:actions}{共轭类}若含于 $A_5$，则要么仍为 $A_5$ 中的一个\omterm{ex:b3:groups:actions}{共轭类}，
要么分裂为两个，取决于某元素在 $S_5$ 中的\omterm{ex:b3:groups:actions}{中心化子}
是否含有奇置换（类的大小 $\abs{\text{ 在 }A_5} =
60/\abs{Z_{A_5}(\sigma)}$，且 $Z_{A_5} = Z_{S_5} \cap A_5$）。
对 $\sigma = (1\,2)(3\,4)$：$\abs{Z_{S_5}(\sigma)} = 120/15 = 8$，
$(1\,2) \in Z_{S_5}(\sigma)$ 为奇置换，故 $\abs{Z_{A_5}} = 4$，
该类含 $60/4 = 15$ 个元素，不分裂。对 $\sigma =
(1\,2\,3)$：$Z_{S_5}(\sigma) \supseteq \langle \sigma \rangle
\times \langle (4\,5)\rangle$，阶为 $6 = 120/20$，故二者相等；
它含奇置换 $(4\,5)$，类的大小为 $60/3 = 20$，不分裂。对
$\sigma$ 为 $5$-轮换：$Z_{S_5}(\sigma) = \langle \sigma\rangle$
（阶为 $120/24 = 5$），全为偶置换：$Z_{A_5}(\sigma) = \langle
\sigma\rangle$，相应的 $A_5$ 中的类含 $60/5 = 12$ 个元素——即
$24$ 个五轮换分裂为\emph{两个}各含 $12$ 个元素的
\omterm{ex:b3:groups:actions}{共轭类}。各类大小依次为 $1, 15, 20, 12, 12$。

(b) \omterm{def:b3:groups:normal}{正规子群} $N$ 是若干\omterm{ex:b3:groups:actions}{共轭类}的并，
其中含 $\{e\}$，且 $\abs N \mid 60$。可能的和 $1 +
(\text{取子集} \{15, 20, 12, 12\})$ 为
\[
1,\ 13,\ 13,\ 16,\ 21,\ 25,\ 28,\ 28,\ 33,\ 36,\ 40,\ 40,\ 45,\
48,\ 48,\ 60 ;
\]
列表中能整除 $60$ 的数只有 $1$ 与 $60$，故 $N =
\{e\}$ 或 $A_5$。

(c) 设 $G$ 为\omterm{def:b3:groups:simple}{单群}，且 $\abs G = 60$。$n_5 \mid 12$，$n_5
\equiv 1 \pmod 5$，故 $n_5 \in \{1, 6\}$。若 $n_5 = 1$，则
西罗 $5$-子群正规，与单性矛盾（$1 < 5 <
60$）。因此 $n_5 = 6$。
\end{solution}

\begin{solution}{exo:b3:groups:12}
(a) 由\omterm{ex:b3:groups:actions}{正规化子}的定义，$P$ 在 $H = N_G(P)$ 中正规，
且它是西罗 $p$-子群（在 $H$ 中：其阶已是
整个 $p$-部分——对 $\abs G$ 如此，对 $\abs H$ 更是如此）。
正规的\omterm{def:b3:groups:sylow}{西罗子群}是唯一的：任一其他\omterm{def:b3:groups:sylow}{西罗子群}
均与之共轭（$H$ 中的西罗第二定理），从而与之相等。

(b) 设 $g \in N_G(H)$。则 $gPg^{-1} \subseteq gHg^{-1} = H$
是 $H$ 的子群，且与 $P$ 同阶：一个西罗
$p$-子群（在 $H$ 中），故由 (a) 知 $gPg^{-1} = P$。于是 $g \in
N_G(P) = H$：即 $N_G(H) \subseteq H$，反向包含
显然。

(c) 设 $H \subseteq N \trianglelefteq G$，且 $N$ 为真子群。
$P$ 是西罗 $p$-子群（在 $N$ 中）；对任意 $g \in G$，
$gPg^{-1} \subseteq N$ 是另一个\omterm{def:b3:groups:sylow}{西罗子群}，故 $gPg^{-1} = nPn^{-1}$，其中
$n \in N$（$N$ 中的西罗第二定理），于是 $n^{-1}g \in N_G(P)
\subseteq N$，从而 $g \in N$，即 $N = G$，矛盾（这就是
\emph{弗拉蒂尼论证}）。对极大子群 $M \supseteq
N_G(P)$：$N_G(M) \supseteq M$ 为 $M$ 或 $G$；若为 $G$，
则 $M \trianglelefteq G$ 为一个真\omterm{def:b3:groups:normal}{正规子群}，
且包含 $N_G(P)$——被上一点排除。故
$N_G(M) = M$。
\end{solution}

\begin{solution}{pb:b3:groups:1}
\textbf{1.} 对 $x, y \in G$：由 $(xy)^2 = e$ 得 $xy = (xy)^{-1} =
y^{-1}x^{-1} = yx$（每个元素都是自己的逆元），故为阿贝尔群。
若将 $G$ 写成加法形式，则它是 $\mathbb
F_2$ 上的向量空间（$2x = 0$，公理即阿贝尔群公理）；若
有限，则它有有限基：$G \cong (\Z/2\Z)^k$，阶为
$2^k$。

\textbf{2.} 阶为 $p^2$ 的群：$G$ 为阿贝尔群
（由\cref{thm:b3:groups:pfixed}）。若某元素的阶为 $p^2$，
则 $G$ 为循环群。否则所有 $x \neq e$ 的阶均为 $p$；加法记号下
$G$ 即为 $\mathbb F_p$ 上的向量空间（$px = 0$），维数为
$2$（共 $p^2$ 个元素）：$G \cong (\Z/p\Z)^2$。

阶为 $8$ 的阿贝尔群，按元素的最大阶 $m$ 分类：$m =
8$：循环群 $\Z/8\Z$。$m = 2$：由第 1 题得 $(\Z/2\Z)^3$。$m = 4$：
设 $x$ 的阶为 $4$，$y \notin \langle x \rangle$；$y^2 \in
\langle x\rangle$（指数为 $2$）。$y^2 \in \{x, x^3\}$ 将使 $y$
的阶为 $8$；故 $y^2 \in \{e, x^2\}$。若 $y^2 = x^2$，则将 $y$ 换为
$xy$：$(xy)^2 = x^2y^2 = x^4 = e$（$G$ 为阿贝尔群），且 $xy \notin
\langle x\rangle$。故可设 $y^2 = e$，于是 $\langle x
\rangle \cap \langle y \rangle = \{e\}$，二者均正规（阿贝尔），
$\abs{\langle x\rangle \langle y\rangle} = 8$：
由\cref{prop:b3:groups:direct} 得 $G \cong \Z/4\Z \times \Z/2\Z$。
两两不同构：方程 $x^2 = e$ 的解数依次为 $2, 4, 8$，
分别对应三个群。

\textbf{3.} 定义 $\psi \colon H \rtimes_{\varphi'} K \to H
\rtimes_{\varphi} K$，由 $\psi(h, k) = (h, \alpha(k))$ 给出，它是双射。
同态性：
\[
\psi\bigl((h,k)(h',k')\bigr)
= \bigl(h\,\varphi'(k)(h'),\, \alpha(kk')\bigr)
= \bigl(h\,\varphi(\alpha k)(h'),\, \alpha(k)\alpha(k')\bigr)
= \psi(h,k)\,\psi(h',k').
\]

\textbf{4.} 由 $\operatorname{Aut}(\Z/n\Z) \cong (\Z/n\Z)^\times$
（见\cref{exo:b3:groups:3}）：当 $n = 3$ 时，$\{\pm 1\} \cong \Z/2\Z$；
当 $n = 4$ 时，$\{\bar 1, \bar 3\} \cong \Z/2\Z$；当 $n = 5$ 时，$\{\bar 1,
\bar 2, \bar 3, \bar 4\}$，为 $4$ 阶循环群，由 $\bar
2$ 生成（$2, 4, 3, 1$）；$n = 7$ 时，为 $6$ 阶循环群，由 $\bar
3$ 生成（$3, 2, 6, 4, 5, 1$）。对 $V = (\Z/2\Z)^2$，其自同构
均为 $\mathbb F_2$-线性的（它保持加法，而标量仅为 $0,
1$），故 $\operatorname{Aut}(V) = GL_2(\mathbb F_2)$，阶为
$(4-1)(4-2) = 6$；它忠实地作用于 $3$ 个非零向量上，
由此得到 $S_3$ 的单同态，两群阶均为
$6$，故 $\operatorname{Aut}(V) \cong S_3$。

\textbf{5.} 由柯西定理给出 $r$，其阶为 $p$；$N = \langle r
\rangle$ 的指数为 $2$，故为\omterm{def:b3:groups:normal}{正规子群}（见\cref{exo:b3:groups:1}）。
又由柯西定理给出 $s$，其阶为 $2$，且 $s \notin N$（
$N$ 的所有非单位元的阶均为奇数 $p$）。

\textbf{6.} $N \cap \langle s \rangle = \{e\}$，且 $\abs{N\langle
s\rangle} = 2p$（见\cref{exo:b3:groups:4}(a)）。由
\cref{prop:b3:groups:semidirect}可知，$G \cong \Z/p\Z \rtimes_\varphi
\Z/2\Z$，其中 $\varphi(1) = (x \mapsto sxs^{-1})$ 为自同构，
其阶整除 $2$。在 $(\Z/p\Z)^\times$ 中，$k^2 = 1$ 的解仅为
$k = \pm 1$（$X^2 - 1$ 在
域 $\mathbb F_p$ 中至多有两个根）。若 $\varphi(1) = \mathrm{id}$，则积
为直积，$G \cong \Z/p\Z \times \Z/2\Z \cong \Z/2p\Z$。若
$\varphi(1) = -\mathrm{id}$: $G = \langle r, s \mid r^p = s^2 = e,
\ srs^{-1} = r^{-1}\rangle \cong D_p$
（见\cref{ex:b3:groups:semidirectexamples}）。两者不同构：
$D_p$ 在 $p \geq 3$ 时非阿贝尔（$srs^{-1} = r^{-1} \neq r$）。

\textbf{7.} $n_q \equiv 1 \pmod q$，它整除 $p < q$：$n_q = 1$，
故 $N \cong \Z/q\Z$ 为\omterm{def:b3:groups:normal}{正规子群}。设 $P \cong \Z/p\Z$ 为一个西罗
$p$-子群：$N \cap P = \{e\}$，$NP = G$（阶为 $pq$），故 $G
\cong \Z/q\Z \rtimes_\varphi \Z/p\Z$，其中 $\varphi \colon \Z/p\Z
\to \operatorname{Aut}(\Z/q\Z) \cong \Z/(q-1)\Z$（循环，在本阶段
不加证明地承认）。若 $p \nmid q - 1$：$\varphi$ 的像之阶
同时整除 $p$ 与 $q-1$，故其平凡，于是 $G \cong
\Z/pq\Z$（见\cref{ex:b3:groups:pq}）。若 $p \mid q - 1$：除
平凡的 $\varphi$ 外，任一非平凡的 $\varphi$ 均为单射（
其核为 $\Z/p\Z$ 的子群，故平凡），其像为\emph{唯一}的
子群 $C$，阶为 $p$，位于循环群
$\Z/(q-1)\Z$ 中。两个非平凡的作用 $\varphi, \varphi'$ 则
为两个同构 $\Z/p\Z \to C$，故 $\alpha =
\varphi^{-1}\circ\varphi' \in \operatorname{Aut}(\Z/p\Z)$
满足 $\varphi' = \varphi\circ\alpha$：由第 3 小题，两个
\omterm{def:b3:groups:semidirect}{半直积}同构。因此恰有一个非阿贝尔
群，阶为 $pq$（非阿贝尔是因为 $\varphi \neq
\mathrm{id}$ 使某个共轭非平凡）。阶为 $15$ 时：$p =
3$，$q = 5$，$3 \nmid 4$：只有循环群。

\textbf{8.} 并非每个元素的阶都 $\leq 2$（否则由
第 1 小题得阿贝尔），且没有元素的阶为 $8$（否则循环，从而阿贝尔）：
存在元素 $r$，其阶为 $4$，且 $N = \langle r\rangle$ 的指数为 $2$，
故为\omterm{def:b3:groups:normal}{正规子群}。

\textbf{9.} 由\cref{exo:b3:groups:1}(b)可知 $s^2 \in N$。
共轭元 $srs^{-1} \in N$ 的阶为 $4$，故 $srs^{-1} \in \{r,
r^3\}$；若 $srs^{-1} = r$，则 $r$ 与 $s$ 可交换，$G =
\langle r, s\rangle$ 为阿贝尔群——排除。故 $srs^{-1} =
r^{-1}$。若 $s^2 = r$ 或 $r^3$，则 $s$ 的阶为 $8$，排除。
（另证：$s^2$ 与 $s$ 可交换，但 $s r s^{-1} =
r^{-1}$，且 $s r^3 s^{-1} = r^{-3} = r$：$r$ 与 $r^3$ 均
不被 $s$ 的共轭作用固定）。故 $s^2 \in \{e, r^2\}$。

\textbf{10.} 若 $s^2 = e$，则 $G = \langle r, s \mid r^4 = s^2 = e,\
srs^{-1} = r^{-1}\rangle$。八个元素 $r^is^j$（$0 \leq i
< 4$，$0 \leq j < 2$）互不相同（$s \notin \langle r\rangle$），
而关系式确定了全部乘积：令 $r
\mapsto$（旋转 $\pi/2$），$s \mapsto$（某个反射）
则给出到 $D_4$ 的满同态，两群阶均为 $8$：
同构。

\textbf{11.} 若 $s^2 = r^2$：同样 $G = \{r^i s^j\}$，并
关系式 $r^4 = e$，$s^2 = r^2$，$srs^{-1} = r^{-1}$ 确定了
整个乘法表。取 $\mathrm i = r$，$\mathrm j = s$，$\mathrm k =
rs$, $-1 = r^2$: $\mathrm i^2 = \mathrm j^2 = -1$, $\mathrm k^2 =
rsrs = r r^{-1} s s = s^2 = -1$（利用 $sr = r^{-1}s$），且
$\mathrm i \mathrm j \mathrm k = r\,s\,rs = r\,r^{-1}s\,s = s^2 =
-1$。存在性：考虑矩阵
\[
A = \begin{pmatrix} \iu & 0\\ 0 & -\iu \end{pmatrix},
\qquad
B = \begin{pmatrix} 0 & 1\\ -1 & 0 \end{pmatrix}
\]
满足 $A^4 = I$，$B^2 = -I = A^2$ 及 $BAB^{-1} = A^{-1}$——
对最后一个关系，验证
\[
BA = \begin{pmatrix} 0 & -\iu\\ -\iu & 0 \end{pmatrix} = A^{-1}B .
\]
于是 $\{\pm I, \pm A, \pm B, \pm AB\}$ 是一个阶为 $8$
的群，实现了该乘法表：$Q_8$ 存在。

\textbf{12.} 设 $H \neq \{e\}$ 为子群，$x \in H
\setminus \{e\}$。若 $x \neq -1$，则 $x \in \{\pm\mathrm i,
\pm\mathrm j, \pm\mathrm k\}$，且 $x^2 = -1 \in H$。故 $-1 \in H$
恒成立。子群为 $\{e\}$、$\{\pm 1\}$（中心）、
$\langle \mathrm i\rangle, \langle \mathrm j\rangle, \langle
\mathrm k\rangle$（指数为 $2$）及 $Q_8$：它们均为\omterm{def:b3:groups:normal}{正规子群}（$\{e\}$ 与
中心显然，指数为 $2$ 者由\cref{exo:b3:groups:1}可知，$Q_8$
本身）。$D_4$ 有五个阶为 $2$ 的元素（$r^2$ 与四个
反射），$Q_8$ 只有一个（$-1$），故不同构。\omterm{def:b3:groups:semidirect}{半直积}
$H \rtimes K$ 中若 $H, K \neq \{e\}$，则须 $H \cap K =
\{e\}$，而两者都含 $-1$，不可能。

\textbf{13.} $n_3 \mid 4$，$n_3 \equiv 1 \pmod 3$：$n_3 \in \{1,
4\}$；$n_2 \mid 3$，为奇数：$n_2 \in \{1, 3\}$。若 $n_3 = 4$，则
四个西罗 $3$-子群两两交集平凡（
阶为素数），给出 $4 \times 2 = 8$ 个阶为 $3$ 的元素；其余
$4$ 个元素必构成唯一的西罗
$2$-子群：$n_2 = 1$。

\textbf{14.} $\ker\rho$ 正规化每个西罗 $3$-子群，故
$\ker \rho \subseteq N_G(P_3)$，后者的指数为 $n_3 = 4$，即
阶为 $3$：$\abs{\ker\rho} \in \{1, 3\}$。若阶为 $3$，则
$\ker\rho$ 将为一个\emph{正规}的西罗 $3$-子群，与 $n_3
= 4$ 矛盾。故 $\rho$ 为单射，其像 $H \leq S_4$ 的阶为
$12$，指数为 $2$：$H \trianglelefteq S_4$，且 $H$ 包含全部
平方元（见\cref{exo:b3:groups:1}(b)）。$S_4$ 的平方元包括
$e$ 与全部八个 $3$-轮换（$\sigma = (\sigma^2)^2$，对一个
$3$-轮换成立），它们生成 $A_4$（它们都在 $A_4$ 中，且
与其乘积一起给出全部十二个元素；或
由\cref{lem:b3:groups:threecycles}（$n = 4$ 时的生成部分，
它只需 $n \geq 3$）。于是 $A_4 \subseteq H$，且 $\abs{A_4}
= \abs H$: $G \cong H = A_4$.

\textbf{15.} $P_3 \trianglelefteq G$, $P_3 \cap P_2 = \{e\}$,
$P_3P_2 = G$: $G \cong \Z/3\Z \rtimes_\varphi P_2$
(\cref{prop:b3:groups:semidirect}), $\varphi \colon P_2 \to
\operatorname{Aut}(\Z/3\Z) = \{\pm\mathrm{id}\} \cong \Z/2\Z$.
\begin{itemize}
  \item $\varphi$ 平凡：直积 $\Z/3\Z \times \Z/4\Z
        \cong \Z/12\Z$ 与 $\Z/3\Z \times (\Z/2\Z)^2 \cong \Z/6\Z
        \times \Z/2\Z$。
  \item $P_2 = \Z/4\Z$，$\varphi$ 为满射：必然
        $\varphi(1) = -\mathrm{id}$（唯一的非平凡选择）：
        一个群，$\mathrm{Dic}_3 = \Z/3\Z \rtimes \Z/4\Z$。
  \item $P_2 = (\Z/2\Z)^2$，$\varphi$ 为满射：$\ker\varphi$ 为
        三个 $2$ 阶子群之一；由此得到的三个
        $\varphi$ 之间只相差 $(\Z/2\Z)^2$ 的
        置换这些子群的自同构（第 4 题：$\operatorname{Aut}
        \cong S_3$ 可迁地作用于三个对合），故
        由第 3 题，它们给出一个同构类。这就是
        $D_6$：取 $t$ 生成 $\ker\varphi$，$x$
        生成 $\Z/3\Z$；元素 $\rho = (x, t)$ 满足
        $\rho^2 = (2x, 0)$，$\rho^3 = (0, t)$，$\rho^6 = e$，且无
        更小的幂为 $e$：阶为 $6$；对 $s = (0, u)$，其中 $u
        \notin \ker\varphi$：$s^2 = e$，且 $s\rho s^{-1} = (-x, t)
        = \rho^{-1}$。由于 $\langle \rho, s\rangle$ 的阶为 $12$，
        故 $G \cong D_6$。
\end{itemize}

\textbf{16.} 计算阶为 $2$ 的元素个数：$\Z/12\Z$ 有 $1$ 个；
$\Z/6\Z\times\Z/2\Z$ 有 $3$ 个；$D_6$ 有 $7$ 个（六个反射与
半转 $\rho^3$）；$A_4$ 有 $3$ 个；$\mathrm{Dic}_3$ 有 $1$ 个
（仅 $(0, 2)$：元素 $(h, k)$ 中 $k$ 为 $4$ 阶元（在
$\Z/4\Z$ 中），其阶为 $4$）。这一计数区分了除
$\{\Z/12\Z, \mathrm{Dic}_3\}$ 与 $\{\Z/6\Z\times\Z/2\Z, A_4\}$ 两对之外的所有群：
每对中前者为交换群，后者不是（$\mathrm{Dic}_3$：
作用非平凡；$A_4$：$(1\,2\,3)$ 与
$(1\,2)(3\,4)$ 不交换）。共五个互异群；
第二至四部分说明此列表是完整的。

\textbf{17.} 分类表：
\begin{center}
\begin{tabular}{c l c}
$n$ & $n$ 阶群 & \#\\
\hline
$1$ & $\{e\}$ & $1$\\
$2$ & $\Z/2\Z$ & $1$\\
$3$ & $\Z/3\Z$ & $1$\\
$4$ & $\Z/4\Z$, $(\Z/2\Z)^2$ & $2$\\
$5$ & $\Z/5\Z$ & $1$\\
$6$ & $\Z/6\Z$, $S_3 = D_3$ & $2$\\
$7$ & $\Z/7\Z$ & $1$\\
$8$ & $\Z/8\Z$, $\Z/4\Z{\times}\Z/2\Z$, $(\Z/2\Z)^3$, $D_4$,
      $Q_8$ & $5$\\
$9$ & $\Z/9\Z$, $(\Z/3\Z)^2$ & $2$\\
$10$ & $\Z/10\Z$, $D_5$ & $2$\\
$11$ & $\Z/11\Z$ & $1$\\
$12$ & $\Z/12\Z$, $\Z/6\Z{\times}\Z/2\Z$, $D_6$, $A_4$,
       $\mathrm{Dic}_3$ & $5$\\
$13$ & $\Z/13\Z$ & $1$\\
$14$ & $\Z/14\Z$, $D_7$ & $2$\\
$15$ & $\Z/15\Z$ & $1$\\
\end{tabular}
\end{center}
阶为 $6, 10, 14$ 的群属第二部分，取 $p = 3, 5, 7$；阶为 $15$ 的
群见第 7 题；阶为 $8$ 的群属第三部分，并见第 2 题；阶为
$12$ 的群属第四部分；素数阶由拉格朗日定理给出；阶为 $4$ 与
$9$ 的群见第 2 题。

\textbf{18.} $Z = Z(G)$ 非平凡
（见\cref{thm:b3:groups:pfixed}）且 $Z \neq G$（非交换），故
$\abs Z \in \{p, p^2\}$。若 $\abs Z = p^2$，则 $G/Z$ 为
$p$ 阶循环群，由\cref{exo:b3:groups:2}知 $G$ 为交换群：
此情形排除，故 $\abs Z = p$，$\abs{G/Z} = p^2$。由第 2 题，
$G/Z$ 为 $\Z/p^2\Z$ 或 $(\Z/p\Z)^2$；循环情形同样
由\cref{exo:b3:groups:2}排除：$G/Z \cong (\Z/p\Z)^2$。由于 $G/Z$
为交换群，故每个\omterm{def:b3:groups:derived}{换位子}都在 $Z$ 中：$D(G) \subseteq Z$；且
$D(G) \neq \{e\}$（$G$ 非交换），故 $D(G) = Z$（$\abs Z = p$
无其他可能）。\omterm{def:b3:groups:derived}{换位子}为中心元，其阶整除
$\abs Z = p$。

\textbf{19.} 对 $k$ 归纳，$k = 1$ 的情形显然。
利用 $yx = xyz^{-1}\cdot$\,——更准确地说，$z = [y, x] =
yxy^{-1}x^{-1}$ 给出 $yx = zxy$，即将一个 $y$ 向左
越过一个 $x$ 便产生一个因子 $z$，它是中心元，
可置于任意位置。于是
\[
(xy)^{k+1} = (xy)^k\,xy = x^ky^kz^{k(k-1)/2}\,xy
= x^k\,(y^kx)\,y\,z^{k(k-1)/2}
= x^{k+1}y^{k+1}\,z^{k(k-1)/2 + k},
\]
由于将 $x$ 移过 $y^k$ 需要 $k$ 个因子 $z$（$y^kx =
z^kxy^k$，由 $k$ 次应用 $yx = zxy$ 得到）；且 $k(k-1)/2 + k
= k(k+1)/2$。

\textbf{20.} 取 $k = p$：$(xy)^p = x^py^pz^{p(p-1)/2}$。对
奇数 $p$，$(p-1)/2$ 为整数，故 $z^{p(p-1)/2} =
(z^p)^{(p-1)/2} = e$（第 18 题：$z$ 的阶整除 $p$）：
$\theta(xy) = \theta(x)\theta(y)$，即一个同态。其值
均为中心元：$x$ 在 $G/Z \cong (\Z/p\Z)^2$ 中的类，其阶整除
$p$，故 $x^p \in Z$。若 $\theta$ 平凡，则每个元素的阶都整除
$p$：方次为 $p$（不为 $1$，因为 $G \neq
\{e\}$）。否则存在 $x^p \neq e$，且 $x$ 的阶为 $p^2$
（阶整除 $p^3$，而 $x$ 不能为 $p^3$ 阶：否则 $G$
将为循环群，从而为交换群）：方次为 $p^2$。

\textbf{21.} 类 $\bar x, \bar y$ 生成 $G/Z$：它们的
\omterm{def:b3:groups:derived}{换位子} $z = [y, x]$ 满足 $\neq e$，否则 $x, y, Z$ 将
生成交换群 $G$（它们的类生成\omterm{thm:b3:groups:quotient}{商群}，
$Z$ 为中心）——且 $z$ 生成 $Z$（$\abs Z = p$）。每个
$g \in G$ 的类可唯一地写为 $\bar x^a\bar y^b$，$0 \leq a, b <
p$，故 $g = x^ay^bz^c$，且 $0 \leq c < p$ 唯一确定：$p^3$
个元素，全部计入。此类标准形之积的计算
仅利用 $yx = zxy$、$z$ 为中心元，以及 $x^p = y^p =
z^p = e$：乘法表唯一确定，故任意两个方次为 $p$
的非交换群（阶为 $p^3$）彼此同构（对应
生成元）。海森堡群实现这些关系：取
$X = I + E_{12}$，$Y = I + E_{23}$，计算得 $[Y, X] = I -
E_{13}$（在 $H_p$ 中为中心元），且对任意严格上
三角矩阵 $N$，由 $N^3 = 0$ 得
\[
(I + N)^p = I + pN + \binom p2N^2 = I
\qquad\text{在特征 } p,\ p \geq 3,
\]
由于 $p \mid p$，且 $p \mid \binom p2$（$p$ 为奇数）：方次
为 $p$。故方次为 $p$ 的群即 $H_p$。

\textbf{22.} $Z(G) = \langle r^p\rangle$：事实上 $r^p$ 是
中心元（练习 20 的论证：$r$ 所在的类在
指数为 $p$ 的\omterm{thm:b3:groups:quotient}{商群} $G/Z$ 中给出 $r^p \in Z$），且 $\neq e$，
因此它生成阶为 $p$ 的中心。任取 $t \notin N$。若
$t^p = e$，则令 $s = t$。否则 $\theta(t) = t^p \in Z =
\langle r^p\rangle$，设 $t^p = r^{pb}$，令 $s = tr^{-b}$，由
练习 20（$\theta$ 为同态，$p$ 为奇数），$s^p =
t^pr^{-pb} = e$，且 $s \notin N$。共轭：$srs^{-1} \in
N$（$N$ 正规）的阶为 $p^2$，故 $srs^{-1} = r^m$，其中 $p
\nmid m$；又 $s^p = e$ 推出 $m^p \equiv m \pmod{p^2}$——
共轭 $p$ 次后回到 $r$，故 $m^p \equiv 1 \pmod
{p^2}$。由费马小定理，$m \equiv m^p \equiv 1 \pmod p$，故 $m = 1 +
ap$。由 $G$ 非交换知 $a \not\equiv 0$；
将 $s$ 换为其幂 $s^{a'}$，其中 $aa' \equiv 1 \pmod p$，
则作用变为 $r \mapsto r^{1+p}$。于是 $G$ 可写成
$\Z/p^2\Z\rtimes_\varphi\Z/p\Z$，其中 $\varphi(1)\colon r
\mapsto r^{1+p}$。由练习 3，任意两个非平凡同态
$\Z/p\Z \to \operatorname{Aut}(\Z/p^2\Z)$，若有相同的像，
——且像恰为\emph{唯一}的阶为 $p$ 的
子群，位于循环群 $\operatorname{Aut}(\Z/p^2\Z)$ 中——
给出同构的\omterm{def:b3:groups:semidirect}{半直积}：唯一性。

\textbf{23.} 交换群：$\Z/p^3\Z$，$\Z/p^2\Z\times\Z/p\Z$，
$(\Z/p\Z)^3$（练习 2 的论证推广一级，按
最大阶分类）。非交换的：恰为 $H_p$（指数为 $p$，
练习 21）和 $\Z/p^2\Z\rtimes\Z/p\Z$（指数为 $p^2$，
练习 22），二者由指数区分。共五个。
当 $p = 2$ 时，练习 20 的同态论证失效：
$z^{p(p-1)/2} = z^{1} = z \neq e$，平方映射不是同态，
事实上，两个 $8$ 阶非交换群的指数都为
$4$——区分 $D_4$ 与 $Q_8$ 的不变量是
阶为 $2$ 的元素个数（五个对一个），而非
指数。这一次，$p$ 为奇数时的情形比
特征 $2$ 的情形更为整齐。

\textbf{24.} 在 $H_p$ 中，每个 $\neq e$ 的元素阶均为 $p$
（指数为 $p$，练习 21）：共 $p^3 - 1$ 个阶为 $p$ 的元素。
在 $M_p$ 中，映射 $\theta \colon x \mapsto x^p$ 是同态
$M_p \to Z(M_p) = \langle r^p\rangle$（练习 20，$p$ 为奇数）；
$\theta(r) = r^p \neq e$，故像为整个阶为 $p$ 的
中心，且 $\ker\theta = \{x : x^p = e\}$ 的阶为 $p^3/p =
p^2$。阶为 $p$ 的元素恰是该核中的
非单位元，共 $p^2 - 1$ 个。当 $p = 3$ 时，$H_3$ 有
$27 - 1 = 26$ 个阶为 $3$ 的元素，$M_3 = \Z/9\Z \rtimes
\Z/3\Z$ 则有 $9 - 1 = 8$ 个。当 $p = 2$ 时，论证
一开始即告失效：非交换的 $8$ 阶群上的平方映射不是同态
（练习 23）。事实上，集合 $\{x : x^2 = e\}$
含 $6$ 个元素（位于 $D_4$ 中）——这一数字不是
$D_4$ 子群的阶。真正能区分这对群的计数是
阶为 $2$ 的元素个数：$D_4$ 中有五个，$Q_8$ 中有一个（练习
11）。

\textbf{25.} 若 $p^2 \mid n$，则群 $\Z/p\Z \times
\Z/(n/p)\Z$ 的阶为 $n$，且不是循环群：每个元素的
阶都整除 $\operatorname{lcm}(p, n/p) = n/p < n$，因为
$p \mid n/p$。若 $p < q$ 为整除 $n$ 的素数，且 $p \mid
q - 1$，则练习 7 给出一个非交换群 $\Z/q\Z \rtimes
\Z/p\Z$，其阶为 $pq$；于是 $(\Z/q\Z \rtimes \Z/p\Z) \times
\Z/(n/pq)\Z$ 的阶为 $n$，且为非交换群，故不是
循环群。现设 $\gcd(n, \varphi(n)) > 1$，取一个素数
$p$，它同时整除两者。记 $\varphi(n) = \prod_{q^a \parallel
n} q^{a-1}(q - 1)$，$p \mid \varphi(n)$ 意味着
要么 $p^2 \mid n$（因子 $q^{a-1}$ 中 $q = p$，
$a \geq 2$），要么 $p \mid q - 1$，其中 $q \mid n$ 是某个素数，且 $q
\neq p$：两种情形下，由上可知 $n$ 均非循环群。
由逆否命题，$n$ 为循环群蕴含 $\gcd(n, \varphi(n)) = 1$。
验证 $n \leq 15$ 的情形：$\varphi(n)$ 对 $n = 1,
\dots, 15$ 的值依次为 $1, 1, 2, 2, 4, 2, 6, 4, 6, 4, 10, 4, 12, 6,
8$，而 $\gcd(n, \varphi(n)) = 1$ 恰对 $n = 1, 2, 3, 5,
7, 11, 13, 15$——恰是练习 17 的表中
仅对应一个群的各项。其余各阶按上述方法
见其非循环：$4, 8, 9, 12$ 由平方因子给出，
$6, 10, 12, 14$ 由 $2 \mid q - 1$ 给出。（逆命题——
$\gcd(n, \varphi(n)) = 1$ 蕴含 $n$ 为循环群——亦成立；
练习 7 证明了其第一个非平凡情形，即 $n = pq$ 且
$p \nmid q - 1$。）
\end{solution}
